% % \smallskip\noindent\textbf{Classical solvers.}
% % The operations research community has developed highly effective solvers over decades, encoding generalizable knowledge in compact, well-understood forms (equations, algorithms, code).
% % Exact and local-search solvers such as LKH-3~\cite{lkhtsp, lkhcvrp} and HGS~\cite{hgs} offer near-optimal precision, while population-based metaheuristics such as Ant Colony Optimization (ACO)~\cite{ACO, maxmin} explore the solution space through pheromone-guided sampling.
% % These methods have been extensively tested in real-world conditions.
% % However, designing and tuning these solvers requires substantial domain expertise, adapting them to new problem variants or dynamic conditions is labor-intensive, and their hand-crafted decision rules can be overly rigid or locally short-sighted.


% \smallskip\noindent\textbf{End-to-end neural solvers.} Typical neural solvers include auto-regressive neural models~\cite{attentionmodel,pomo} and non-auto-regressive neural models~\cite{bq,lehd,invit,sigd,l2c}, which learn constructive policies directly from data, automating solution generation and enabling batch optimization. In principle, NCO can optimize with minimal manual re-engineering, since the learned policy adapts to new instances without re-designing heuristics. In practice, however, they suffer from optimality gaps on hard instances, difficulty generalizing across problem sizes and distributions, and inefficiency on large-scale instances. For example, scaling beyond a few thousand nodes remains challenging due to quadratic attention complexity and error accumulation along long construction trajectories (Section~\ref{sec:related_scalability}).

% \smallskip\noindent\textbf{Learning-guided optimization.}
% These methods~\cite{learn2branch, learn2delegate, neurolkh} combine the strengths of both paradigms: the convergence guarantees and problem-specific structure of classical solvers, together with the data-driven adaptability of neural methods.
% Rather than replacing the solver, a neural network supplies signals while the solver retains its algorithmic backbone, mitigating the scalability limitations of end-to-end NCO and the rigidity of hand-crafted rules. 
% % ACO's factorized transition rule, provides a particularly clean interface for injecting learned guidance, making it a natural bridge to learning-guided optimization.

% \smallskip\noindent\textbf{Neural Guided ACO.}  Ant Colony Optimization (ACO)~\cite{ACO, maxmin} is a population-based metaheuristic that explores solution space through pheromone-guided sampling. Recent works integrate neural components into ACO through diverse mechanisms: DeepACO \cite{deepaco} and GTG-ACO \cite{gtgaco} employ GNNs and Transformers to learn initialization of the heuristic matrices and/or the pheromone matrices, GFACS \cite{gfacs} leverages GFlowNets for constructive sampling, and HeatACO \cite{heataco} utilizes pre-computed heatmaps as probabilistic priors. Despite these architectural distinctions, existing methods suffer from two primary limitations: (1) they rely on static guidance that remains invariant to the search trajectory, and (2) their scalability to large-scale instances is still restricted.
% % Within ACO, neural-guided variants \cite{deepaco, gfacs, gtgaco, heataco} train a model to emit edge heuristics, but all are \emph{static} guidance.
% % Recent finding ~\cite{heataco} shows that the model architecture is interchangeable under this paradigm, suggesting that the static integration, rather than the model, is the limiting factor.
% % Training on post-LS rewards further compounds difficulties; DeepACO's Neural Local Search mitigates gradient washout but is impractical beyond a few thousand nodes.
% % \DyNACO{}'s trajectory-aware objective targets both the static-guidance and credit-assignment limitations.

% % \smallskip\noindent\textbf{Scaling combinatorial optimization.}
% % Decomposition~\cite{htsp,glop} and hierarchical construction~\cite{sil,elg} tackle large instances, while perturbation-based ACO~\cite{partialACO, faco, P-ACO} improve ACO scalability through sparse candidate graphs.
% % \DyNACO{} conditions on per-edge pheromone features and incumbent structure within a perturbation-based backend, enabling fine-grained, learned guidance at scales up to 100K nodes.

\smallskip\noindent\textbf{End-to-end neural solvers.}
End-to-end neural combinatorial optimization methods learn constructive policies directly from data, aiming to automate solution generation with minimal manual re-engineering. Representative approaches include auto-regressive models, which sequentially select nodes or edges~\cite{point,attentionmodel,pomo,mdam,symnco,bq,lehd,l2c}, and non-auto-regressive models, which predict edge-level scores or heatmaps in parallel~\cite{gcn,dimes,difusco,t2t,fastt2t}. Recent works further improve scalability and generalization through heavy decoders~\cite{bq,lehd,reld}, distance-aware or local-view modeling~\cite{elg,invit,icam,dgl}, self-improvement~\cite{sil}, and decomposition-based solving~\cite{htsp,glop,udc}. In principle, these methods can learn implicit heuristic rules that transfer across instances without manually redesigning solvers. In practice, however, their zero-shot generalization remains limited: models trained on narrow synthetic distributions can degrade substantially on benchmark instances with different scales or structures, and large-scale inference is often constrained by attention complexity, sequential error accumulation, or expensive iterative reconstruction. These limitations motivate hybrid designs that preserve the robustness of classical search while using learning to improve selected decision rules.

\smallskip\noindent\textbf{Learning-guided optimization.}
Learning-guided optimization methods~\cite{learn2branch,learn2delegate,neurolkh,rgls} combine the strengths of classical algorithms and neural models. Rather than replacing the entire solver, they use neural networks to provide auxiliary signals, such as branching scores, decomposition decisions, candidate edges, or local-search priorities, while retaining the algorithmic backbone of an established optimizer. Feasibility handling, iterative improvement, and convergence behavior remain governed by the underlying heuristic, which improves scalability and reliability relative to fully neural solvers. From the viewpoint of neural routing solvers, such methods replace selected handcrafted decision rules in classical heuristic frameworks with learned rules while preserving the search structure of the original algorithms. Learning-guided optimization is therefore especially attractive for zero-shot deployment, where a learned component must remain useful under distribution shift instead of fitting a fixed synthetic training distribution.

\smallskip\noindent\textbf{Neural-guided ACO.}
Ant Colony Optimization (ACO)~\cite{ACO,maxmin} is a population-based metaheuristic that explores the solution space through pheromone-guided sampling. Its transition rule naturally factorizes into pheromone information, heuristic priors, and problem-specific feasibility constraints, providing a clean interface for neural guidance. Recent neural-guided ACO methods integrate learning into this interface through different mechanisms. DeepACO~\cite{deepaco} learns heuristic measures with graph neural networks, GFACS~\cite{gfacs} leverages GFlowNets to guide ant sampling, GTG-ACO~\cite{gtgaco} learns heuristic and pheromone-related matrices with Transformer-based modeling, and HeatACO~\cite{heataco} studies heatmap-based priors for ACO search. Despite their architectural differences, these methods typically produce edge-level guidance from the input instance only once, before the main ACO loop starts. The learned signal is therefore static, while the actual search process is dynamic: pheromones evolve, incumbents change, and the population may shift from exploration to exploitation or stagnate around local optima. This creates a training--inference misalignment and can weaken generalization, since a heatmap learned from one distribution may not remain reliable across the diverse search trajectories induced by out-of-distribution instances.

\smallskip\noindent\textbf{Diversity and trajectory-level search.}
Diversity is a central mechanism in population-based heuristics, where multiple candidate solutions are maintained to balance intensification and diversification. In neural combinatorial optimization, early diversity mechanisms often exploit symmetries of the solution representation. POMO samples multiple rollouts by fixing different first actions~\cite{pomo}, which is effective for symmetric routing problems but couples diversity to the construction order. Later methods make diversity more policy-driven. MDAM uses multiple decoders and divergence-based training to encourage different construction distributions~\cite{mdam}. Poppy trains a population of policies by updating the member that performs best on each instance, so specialization is induced by population performance rather than by a predefined diversity metric~\cite{poppy}. PolyNet further reduces the cost of learning multiple strategies by using a single decoder conditioned on strategy identifiers, showing that complementary solution strategies can emerge without forcing diverse first moves~\cite{polynet}.

These methods provide important background for learned strategy diversity. In construction policies, diversity appears mainly as a set of complete solutions sampled from different neural distributions. In neural-guided ACO, strategy diversity can instead be expressed through behavior-conditioned transition priors used by ant groups within one shared colony. The ant groups can follow different state-conditioned priors while continuing to cooperate through the same incumbent solution, pheromone update, and local refinement mechanism.

\smallskip\noindent\textbf{Scaling combinatorial optimization.}
Scaling neural solvers to large routing instances remains challenging. 
End-to-end constructive policies~\cite{attentionmodel,pomo,bq,lehd, invit, sigd} are typically trained on small-scale instances ($N{=}50$ or $100$) and rely on quadratic attention complexity $O(N^2)$ in the encoder.
When applied to large-scale instances ($N{\geq}1{,}000$), these models suffer from out-of-distribution degradation, as the learned representations do not transfer across scales, compounded by memory exhaustion from the $O(N^2)$ encoder. Recent methods address this through divide-and-conquer (GLOP~\cite{glop}, H-TSP~\cite{htsp}) or hierarchical construction (SIL~\cite{sil}, LEHD~\cite{lehd}).
However, decomposition methods depend on global context embeddings that degrade at extreme scales, and hierarchical methods incur prohibitive runtimes (SIL PRC$_{1000}$: 42h at TSP-100K).
% End-to-end approaches often rely on decomposition, search-space reduction, or hierarchical construction~\cite{htsp,glop,udc,sil,elg,invit,l2r,dgl},
In classical ACO, its variants often improve scalability through candidate lists, sparse graphs, and perturbation-based updates~\cite{partialACO,faco,PACO}. However, in learning-guided ACO, scalability also depends on whether post-refinement rewards remain attributable to the neural decisions that produced them, which make it harder to scale to large instances without sacrificing local search convergent.
% \DyNACO{} combines sparse candidate graphs, perturbation-based ant sampling, and scope-restricted refinement so that the policy controls localized solution changes with meaningful post-refinement feedback, enabling dynamic neural guidance at scales up to 100K nodes.
